% ============================================================================
% templates/exam/examen-objetivo/examen-objetivo.tex — plantilla: examen objetivo
% ----------------------------------------------------------------------------
% Para: opción múltiple, verdadero/falso, emparejamiento, completar espacios,
% ordenamiento y selección múltiple.
% Uso: scripts/build.sh examen-objetivo.tex   (LuaLaTeX)
% Con claves: \documentclass[claves]{academic-exam}
% ============================================================================
\documentclass{academic-exam}

\universidad{Universidad Nacional de San Cristóbal de Huamanga}
\facultad{Facultad de Ciencias Económicas, Administrativas y Contables}
\escuela{Escuela Profesional de Economía}
\curso{Microeconomía I}
\codigocurso{EC-241}
\docente{Mg. Nombre Apellido}
\periodo{Semestre 2026-I}
\tipoevaluacion{Examen Virtual 01}
\fechaevaluacion{21 de julio de 2026}
\duracion{60 minutos}
\modalidad{Virtual · sincrónico}

\begin{document}
\membrete

\begin{instrucciones}
\begin{itemize}
  \item Marque con un aspa (\texttimes) o rellene la casilla correspondiente.
        Solo hay una alternativa correcta por pregunta, salvo indicación
        contraria.
  \item Las respuestas con borrones o doble marca se consideran nulas.
  \item Puntaje total: \totalpuntos{} puntos en \totalpreguntas{} preguntas.
        No hay puntos en contra.
\end{itemize}
\end{instrucciones}

% ================================================== I. OPCIÓN MÚLTIPLE =====
\section*{\color{colorprimario}Parte I · Opción múltiple}

\begin{pregunta}[2]
Si el precio de un bien Giffen aumenta, la cantidad demandada del bien:
\begin{opciones}
  \opcion disminuye, como predice la ley de la demanda.
  \opcion \correcta{aumenta, porque el efecto ingreso domina al efecto
          sustitución.}
  \opcion no cambia, porque la demanda es perfectamente inelástica.
  \opcion aumenta solo si el bien es de lujo.
\end{opciones}
\clave{b}
\end{pregunta}

\begin{pregunta}[2]
La elasticidad precio de la demanda en el punto medio de una curva de
demanda lineal es:
\begin{opciones}(4)
  \opcion $0$
  \opcion $-0{,}5$
  \opcion \correcta{$-1$}
  \opcion $-\infty$
\end{opciones}
\clave{c}
\end{pregunta}

\begin{pregunta}[2]
Seleccione \textbf{todas} las afirmaciones correctas sobre un mercado en
competencia perfecta (selección múltiple):
\begin{opciones}(1)
  \opcion \correcta{Las empresas son precio-aceptantes.}
  \opcion Existen barreras de entrada significativas.
  \opcion \correcta{El producto es homogéneo.}
  \opcion Cada empresa enfrenta una demanda de pendiente negativa.
\end{opciones}
\clave{a, c}
\end{pregunta}

% ================================================ II. VERDADERO / FALSO ====
\section*{\color{colorprimario}Parte II · Verdadero o falso}

\begin{pregunta}[4]
Indique si cada enunciado es verdadero (V) o falso (F)
\emph{(1 punto por acierto)}:
\begin{tablavf}
  \vf[F]{El excedente del consumidor aumenta cuando sube el precio de
         mercado.}
  \vf[V]{En el óptimo del consumidor, la relación marginal de sustitución
         iguala la razón de precios.}
  \vf[V]{Un impuesto específico genera pérdida irrecuperable de eficiencia
         salvo que la demanda o la oferta sean perfectamente inelásticas.}
  \vf[F]{Los bienes inferiores tienen elasticidad ingreso positiva.}
\end{tablavf}
\end{pregunta}

% ================================================= III. EMPAREJAMIENTO ====
\section*{\color{colorprimario}Parte III · Emparejamiento}

\begin{pregunta}[4]
Relacione cada concepto de la columna A con su definición en la columna B,
escribiendo la letra correspondiente entre los paréntesis:
\begin{emparejamiento}
  \pareja[c]{Costo de oportunidad}
    {Curva que muestra las combinaciones de dos bienes que reportan la
     misma utilidad.}
  \pareja[a]{Curva de indiferencia}
    {Valor de la mejor alternativa a la que se renuncia al tomar una
     decisión.}
  \pareja[d]{Restricción presupuestaria}
    {Cantidad adicional de utilidad que aporta la última unidad consumida.}
  \pareja[b]{Utilidad marginal}
    {Conjunto de canastas que el consumidor puede adquirir con su ingreso.}
\end{emparejamiento}
\end{pregunta}

% ============================================ IV. COMPLETAR Y ORDENAR ======
\section*{\color{colorprimario}Parte IV · Completar y ordenar}

\begin{pregunta}[3]
Complete los espacios en blanco:
\begin{apartados}
  \item Cuando la elasticidad precio de la demanda es mayor que uno en
        valor absoluto, se dice que la demanda es
        \completar[elástica]{3.2cm} y una reducción del precio
        \completar[aumenta]{3.2cm} el ingreso total del vendedor.
  \item El punto donde la curva de costo marginal corta al costo medio
        corresponde al \completar[mínimo del costo medio]{4.5cm}.
\end{apartados}
\end{pregunta}

\begin{pregunta}[3]
Ordene cronológicamente las etapas del análisis de equilibrio parcial
escribiendo los números 1 a 4:
\begin{apartados}
  \item (\,\completar[3]{0.9cm}\,) Determinación del precio de equilibrio.
  \item (\,\completar[1]{0.9cm}\,) Especificación de la función de demanda.
  \item (\,\completar[2]{0.9cm}\,) Especificación de la función de oferta.
  \item (\,\completar[4]{0.9cm}\,) Análisis de estática comparativa.
\end{apartados}
\end{pregunta}

\findelexamen
\end{document}
